% ID: MAT-TRI-IDN-001
% Titolo: Identità con seno coseno e tangente
% Disciplina: Matematica
% Area: Trigonometria
% Argomento: Identita
% Classe: Quarta secondaria di secondo grado
% Difficolta: 3
% Tipo: dimostrazione
% Tag: matematica, trigonometria, identita
% Fonte: verifica_funzioni_goniometriche_anonimizzata
% Provenienza: materiale fornito dal docente; testo anonimizzato; soluzione redatta per Exergo

\begin{esercizio}
Verifica l'identit\`a
    \[
    \frac{\sin\alpha\cos\alpha}{2}
    +2\tan\alpha\cos^2\alpha
    -(\sin\alpha+\cos\alpha)^2
    =
    \frac{\tan\alpha}{2}\left(1-\sin^2\alpha\right)-1.
    \]
\end{esercizio}

\begin{soluzione}
Il dominio comune impone $\cos\alpha\ne0$.
Il primo membro diventa $\tfrac12\sin\alpha\cos\alpha+2\sin\alpha\cos\alpha
-1-2\sin\alpha\cos\alpha=\tfrac12\sin\alpha\cos\alpha-1$.
Il secondo è $\tfrac12(\sin\alpha/\cos\alpha)\cos^2\alpha-1$, lo stesso valore.
\end{soluzione}
